\chapter{Linux BSP 与设备驱动}
\label{chap:linux-bsp}

板级支持包（Board Support Package, BSP）负责把处理器、SoC 外设和电路板资源接入 Bootloader、Linux 内核及用户空间。
BSP 不是若干补丁的集合，而是一组可维护的硬件描述、驱动、配置、固件和验证资产。
本章讨论主线 Linux 常见模型；具体数据结构和 API 会随内核版本演进，应以目标源码树中的文档为准\cite{linux-driver-api}。

\section{BSP 的边界与交付物}

一个可交付 BSP 通常包含：
\begin{itemize}
  \item Bootloader 板级配置、存储布局和启动策略；
  \item Linux defconfig、设备树、必要驱动和固件；
  \item 时钟、复位、供电、引脚、中断和 DMA 等资源描述；
  \item 根文件系统中的固件、规则、服务和诊断工具；
  \item 可复现的构建配置、镜像清单和版本信息；
  \item 上电、接口、压力、低功耗和恢复测试结果。
\end{itemize}

BSP 应尽量复用上游内核接口。长期依赖大规模 vendor fork 会增加安全更新、内核升级和问题定位成本。

\section{Linux 设备模型}

Linux 设备模型以设备（\texttt{struct device}）、驱动（\texttt{struct device\_driver}）、总线和类为核心，统一管理设备发现、匹配、生命周期、电源管理和用户空间表示。

\subsection{设备与驱动匹配}

设备可以由可枚举总线发现，也可以由固件描述创建。SoC 内部不可枚举外设通常注册为 platform device，并通过 Device Tree 的 \texttt{compatible} 与驱动匹配。
\texttt{compatible} 描述硬件编程模型，而不是板卡营销名称；兼容字符串应从具体到通用排列，以便驱动选择精确实现或兼容回退。

驱动的 \texttt{probe()} 应完成资源获取、硬件初始化和子系统注册。
现代驱动优先使用设备管理资源接口，使时钟、映射、IRQ、内存和注册对象在 probe 失败或设备解绑时按正确顺序释放。

\subsection{Deferred Probe}

设备之间存在供电、时钟、PHY、IOMMU 等依赖。如果消费者 probe 时供应者尚未就绪，资源获取函数可能返回 \texttt{-EPROBE\_DEFER}。
驱动应传播该错误，而不是把暂时缺失误判为永久硬件故障。反复 deferred probe 通常意味着依赖描述、驱动配置或供应者初始化存在问题。

\section{Device Tree 工程}

Device Tree 描述硬件拓扑和操作系统发现硬件所需的信息，不用于表达运行策略或复制驱动内部默认值。
语法和通用语义由 Devicetree Specification 定义\cite{devicetree-spec}，Linux 绑定还应遵守目标内核的 YAML schema。

\subsection{节点与资源}

典型外设节点包含地址、中断、时钟、复位、DMA、供电和引脚资源：

\begin{lstlisting}[language=dts]
serial@40010000 {
    compatible = "vendor,soc-uart-v2", "vendor,soc-uart";
    reg = <0x40010000 0x1000>;
    interrupts = <42>;
    clocks = <&clk UART0_CLK>;
    resets = <&rst UART0_RST>;
    pinctrl-names = "default", "sleep";
    pinctrl-0 = <&uart0_default>;
    pinctrl-1 = <&uart0_sleep>;
    status = "okay";
};
\end{lstlisting}

单元地址应与 \texttt{reg} 的首地址一致；\texttt{status}、父总线的地址单元和范围转换必须符合层次关系。
设备树变更应通过 \texttt{dtbs\_check} 和目标板启动验证，不能只以 DTS 能够编译作为完成标准。

\subsection{绑定设计原则}

新绑定应描述稳定的硬件接口。能够由寄存器或标准机制发现的信息通常不应重复写入设备树；板级连线、不可发现参数和真实硬件差异则应显式描述。
绑定一旦进入公开 ABI，就必须考虑旧设备树与新内核、新设备树与旧内核的兼容性。

\section{SoC 基础资源}

\subsection{引脚、时钟、复位与供电}

pinctrl 子系统管理复用、电气属性和功耗状态；common clock framework 管理时钟拓扑、频率和门控；reset controller 表达复位线；regulator framework 管理电源轨、电压和消费者关系。

驱动不应绕过这些框架直接操作其他模块的寄存器。推荐的启用顺序通常是获取资源、设置所需速率或电压、解除复位、使能时钟并初始化设备；关闭路径必须按硬件约束逆序执行。

\subsection{中断域}

Linux 使用 \texttt{irq\_domain} 把硬件中断号映射为内核虚拟 IRQ。级联控制器和 GPIO 中断控制器通常建立层次中断域。
驱动只应使用固件接口和 IRQ API 获得 IRQ，不应假定 Linux IRQ 号等于硬件中断号。

中断处理应根据硬件特性选择硬中断、线程化中断或工作队列。硬中断部分负责确认设备并保存最少状态；可能睡眠的总线访问或复杂处理应放入线程上下文。

硬件中断号（hwirq）只在某个中断控制器内部有意义，Linux IRQ 则是内核分配的逻辑编号。\texttt{irq\_domain} 负责二者之间的映射，并把 Device Tree 或 ACPI 中的中断描述转换为控制器资源。GPIO 控制器、PMIC 和外部 FPGA 经常作为级联 irqchip，其 domain 是上级 GIC domain 的子节点。

固定且较小的 hwirq 空间通常使用线性映射；稀疏而巨大的编号空间可以使用树形映射；层次中断控制器应通过 \texttt{irq\_domain\_create\_hierarchy()} 建立父子关系。驱动不得缓存或导出动态 Linux IRQ 编号，调试时应同时记录 domain、hwirq、触发类型和目标 CPU。

\begin{lstlisting}[language=C,caption={级联中断控制器的 domain 建立骨架}]
static const struct irq_domain_ops child_domain_ops = {
    .alloc = child_irq_alloc,
    .free = irq_domain_free_irqs_common,
    .translate = irq_domain_translate_twocell,
};

static int child_irqchip_probe(struct platform_device *pdev)
{
    struct irq_domain *parent, *domain;
    struct device_node *parent_np;

    parent_np = of_irq_find_parent(pdev->dev.of_node);
    if (!parent_np)
        return -ENXIO;
    parent = irq_find_host(parent_np);
    of_node_put(parent_np);
    if (!parent)
        return -EPROBE_DEFER;

    domain = irq_domain_create_hierarchy(parent, 0, 32,
        dev_fwnode(&pdev->dev), &child_domain_ops,
        platform_get_drvdata(pdev));
    return domain ? 0 : -ENOMEM;
}
\end{lstlisting}

\subsection{DMA 与 IOMMU}

设备看到的 DMA 地址不一定等于 CPU 物理地址。DMA API 负责缓存一致性、地址转换、bounce buffer 和 IOMMU 映射，驱动不能直接用 \texttt{virt\_to\_phys()} 替代 DMA 映射。

一致性缓冲区适合长期由 CPU 和设备共享的控制结构；流式映射适合具有明确所有权切换的数据缓冲区。
调用同步 API 后，驱动必须遵守 CPU 与设备之间的所有权协议。

scatter-gather 列表用于描述非连续内存。驱动应使用 DMA 映射后返回的段数进行硬件编程，因为映射层可能合并相邻条目。

\subsection{Scatter-gather DMA 的所有权}

\texttt{struct scatterlist} 描述页、页内偏移和长度；\texttt{struct sg\_table} 同时保存原始条目数 \texttt{orig\_nents} 与 DMA 映射后的条目数 \texttt{nents}。驱动调用 \texttt{dma\_map\_sg()} 后，必须使用返回值或 \texttt{nents} 对硬件描述符编程，而不能继续使用原始段数。

\begin{lstlisting}[language=C,caption={流式 scatter-gather DMA 的标准生命周期}]
int mapped = dma_map_sg(dev, sgl, orig_nents, DMA_TO_DEVICE);
if (mapped == 0)
    return -EIO;

for_each_sg(sgl, sg, mapped, i)
    program_desc(i, sg_dma_address(sg), sg_dma_len(sg));

start_dma();
wait_for_completion(&done);
dma_unmap_sg(dev, sgl, orig_nents, DMA_TO_DEVICE);
\end{lstlisting}

映射期间缓冲区归设备所有，CPU 不应访问；重新交给 CPU 前，应完成传输并按 DMA API 要求同步或解除映射。错误路径、超时路径和设备移除路径必须保证只 unmap 一次，并在释放页之前终止 DMA。

\section{驱动执行上下文}

Linux 内核代码可能运行在进程上下文、硬中断、软中断、工作队列或专用内核线程中。
能否睡眠、允许使用的锁以及可调用 API 取决于当前上下文。

\subsection{锁与并发}

互斥锁适用于可以睡眠的进程上下文；自旋锁适用于短小且不能睡眠的临界区；原子变量和位操作只解决有限的状态更新问题，不能自动建立完整对象生命周期。

驱动必须考虑以下并发来源：
\begin{itemize}
  \item 多个用户线程同时调用文件操作；
  \item 中断与进程上下文访问共享状态；
  \item DMA 完成、超时和设备移除同时发生；
  \item runtime PM 与正常 I/O 竞争；
  \item CPU 热插拔、系统挂起和驱动解绑。
\end{itemize}

锁顺序和对象所有权应写入设计说明。不要通过大范围关中断或单个全局锁掩盖生命周期问题。

\subsection{等待与异步工作}

等待队列用于等待条件变化；completion 适合一次性或阶段性完成事件；workqueue 用于把工作推迟到可睡眠上下文；timer 和 hrtimer 用于时间触发。
所有异步工作在设备移除和错误回滚路径中都必须被取消或同步，避免释放后的回调访问失效对象。

\section{常见子系统接入}

驱动应优先接入语义匹配的内核子系统，而不是为每个设备创建私有字符设备接口：
\begin{itemize}
  \item 传感器和 ADC：IIO；
  \item 按键、触摸和输入设备：input；
  \item 温度、电压和风扇监控：hwmon；
  \item 非易失存储：MTD、SPI NOR、MMC；
  \item 网络设备：netdev、PHY 和 phylink；
  \item 摄像头：V4L2/media；
  \item 音频：ALSA SoC；
  \item 显示：DRM/KMS；
  \item 看门狗、RTC、PWM、LED：对应标准子系统。
\end{itemize}

使用标准子系统能够复用公共 ABI、电源管理、诊断工具和测试基础设施，也降低应用与单一驱动实现的耦合。

\section{完整 platform driver 骨架}

下面的驱动展示一个由 Device Tree 枚举、具有 MMIO、时钟、复位和线程化中断的 platform device。示例省略具体寄存器语义，但保留 probe 回滚、runtime PM、IRQ 与 remove 生命周期，能够作为接入 IIO、hwmon、input 等标准子系统之前的结构模板。

\begin{lstlisting}[language=C,caption={带 runtime PM 的 platform driver 结构}]
#include <linux/clk.h>
#include <linux/interrupt.h>
#include <linux/io.h>
#include <linux/module.h>
#include <linux/of.h>
#include <linux/platform_device.h>
#include <linux/pm_runtime.h>
#include <linux/reset.h>

struct sample_dev {
    struct device *dev;
    void __iomem *base;
    struct clk *clk;
    struct reset_control *rst;
    int irq;
};

static irqreturn_t sample_irq(int irq, void *data)
{
    struct sample_dev *sdev = data;
    u32 status = readl(sdev->base + STATUS_REG);

    if (!(status & IRQ_PENDING))
        return IRQ_NONE;
    writel(status, sdev->base + STATUS_REG); /* acknowledge */
    return IRQ_WAKE_THREAD;
}

static irqreturn_t sample_irq_thread(int irq, void *data)
{
    struct sample_dev *sdev = data;

    sample_push_completed_data(sdev); /* may sleep */
    return IRQ_HANDLED;
}

static int sample_runtime_resume(struct device *dev)
{
    struct sample_dev *sdev = dev_get_drvdata(dev);
    int ret;

    ret = clk_prepare_enable(sdev->clk);
    if (ret)
        return ret;
    ret = reset_control_deassert(sdev->rst);
    if (ret)
        clk_disable_unprepare(sdev->clk);
    return ret;
}

static int sample_runtime_suspend(struct device *dev)
{
    struct sample_dev *sdev = dev_get_drvdata(dev);

    reset_control_assert(sdev->rst);
    clk_disable_unprepare(sdev->clk);
    return 0;
}

static int sample_probe(struct platform_device *pdev)
{
    struct device *dev = &pdev->dev;
    struct sample_dev *sdev;
    int ret;

    sdev = devm_kzalloc(dev, sizeof(*sdev), GFP_KERNEL);
    if (!sdev)
        return -ENOMEM;
    sdev->dev = dev;
    sdev->base = devm_platform_ioremap_resource(pdev, 0);
    if (IS_ERR(sdev->base))
        return PTR_ERR(sdev->base);
    sdev->clk = devm_clk_get(dev, NULL);
    if (IS_ERR(sdev->clk))
        return dev_err_probe(dev, PTR_ERR(sdev->clk), "clock\n");
    sdev->rst = devm_reset_control_get_optional_exclusive(dev, NULL);
    if (IS_ERR(sdev->rst))
        return dev_err_probe(dev, PTR_ERR(sdev->rst), "reset\n");
    sdev->irq = platform_get_irq(pdev, 0);
    if (sdev->irq < 0)
        return sdev->irq;

    platform_set_drvdata(pdev, sdev);
    ret = devm_request_threaded_irq(dev, sdev->irq, sample_irq,
        sample_irq_thread, IRQF_ONESHOT, dev_name(dev), sdev);
    if (ret)
        return dev_err_probe(dev, ret, "irq\n");

    pm_runtime_set_suspended(dev);
    pm_runtime_enable(dev);
    pm_runtime_set_autosuspend_delay(dev, 100);
    pm_runtime_use_autosuspend(dev);
    ret = pm_runtime_resume_and_get(dev);
    if (ret < 0) {
        pm_runtime_disable(dev);
        return ret;
    }
    ret = sample_register_subsystem(sdev);
    if (ret) {
        pm_runtime_put_sync_suspend(dev);
        pm_runtime_disable(dev);
        return ret;
    }
    pm_runtime_mark_last_busy(dev);
    pm_runtime_put_autosuspend(dev);
    return 0;
}

static void sample_remove(struct platform_device *pdev)
{
    struct sample_dev *sdev = platform_get_drvdata(pdev);
    int ret;

    ret = pm_runtime_resume_and_get(&pdev->dev);
    sample_unregister_subsystem(sdev);
    if (ret >= 0)
        pm_runtime_put_sync_suspend(&pdev->dev);
    pm_runtime_disable(&pdev->dev);
}

static const struct dev_pm_ops sample_pm_ops = {
    SET_RUNTIME_PM_OPS(sample_runtime_suspend,
                       sample_runtime_resume, NULL)
};

static const struct of_device_id sample_of_match[] = {
    { .compatible = "vendor,sample-v1" }, { }
};
MODULE_DEVICE_TABLE(of, sample_of_match);

static struct platform_driver sample_driver = {
    .probe = sample_probe,
    .remove_new = sample_remove,
    .driver = {
        .name = "sample-device",
        .of_match_table = sample_of_match,
        .pm = &sample_pm_ops,
    },
};
module_platform_driver(sample_driver);
MODULE_LICENSE("GPL");
\end{lstlisting}

真实驱动还必须在注册子系统之前确认 runtime PM 初始状态一致，并根据目标内核版本选择资源获取 API。probe 的每一个失败点、模块卸载、设备解绑、系统休眠和 IRQ 并发都应进入测试矩阵。

\section{标准子系统驱动实例：IIO ADC}

ADC 不应默认设计私有字符设备。IIO 能统一表达通道、raw 值、scale、采样频率、触发和缓冲采集。下面给出一个 MMIO ADC 的最小接入结构，寄存器值仅用于说明接口关系。

设备树绑定应先定义硬件资源：

\begin{lstlisting}[language=dts]
adc@48000000 {
    compatible = "vendor,soc-adc-v1";
    reg = <0x48000000 0x1000>;
    interrupts = <52>;
    clocks = <&clk 12>;
    vref-supply = <&vref_1v8>;
    #io-channel-cells = <1>;
};
\end{lstlisting}

驱动把硬件通道转换为 IIO channel，并通过 \texttt{read\_raw()} 暴露标准属性：

\begin{lstlisting}[language=C,caption={IIO ADC 的核心接口}]
static const struct iio_chan_spec adc_channels[] = {
    { .type = IIO_VOLTAGE, .indexed = 1, .channel = 0,
      .info_mask_separate = BIT(IIO_CHAN_INFO_RAW),
      .info_mask_shared_by_type = BIT(IIO_CHAN_INFO_SCALE) },
    { .type = IIO_VOLTAGE, .indexed = 1, .channel = 1,
      .info_mask_separate = BIT(IIO_CHAN_INFO_RAW),
      .info_mask_shared_by_type = BIT(IIO_CHAN_INFO_SCALE) },
};

static int adc_read_raw(struct iio_dev *indio,
                        const struct iio_chan_spec *chan,
                        int *val, int *val2, long mask)
{
    struct adc_state *st = iio_priv(indio);
    int ret;

    switch (mask) {
    case IIO_CHAN_INFO_RAW:
        ret = pm_runtime_resume_and_get(st->dev);
        if (ret < 0)
            return ret;
        ret = adc_single_conversion(st, chan->channel, val);
        pm_runtime_mark_last_busy(st->dev);
        pm_runtime_put_autosuspend(st->dev);
        return ret ? ret : IIO_VAL_INT;
    case IIO_CHAN_INFO_SCALE:
        *val = st->vref_uv;
        *val2 = st->resolution_bits;
        return IIO_VAL_FRACTIONAL_LOG2;
    default:
        return -EINVAL;
    }
}

static const struct iio_info adc_info = {
    .read_raw = adc_read_raw,
};
\end{lstlisting}

probe 需要分配 \texttt{iio\_dev}、映射资源、获取参考电压和时钟、初始化 completion/IRQ、设置通道与 \texttt{iio\_info}，最后调用 \texttt{devm\_iio\_device\_register()}。单次转换应在启动硬件前重新初始化 completion，等待时使用有界超时，并在超时后停止转换和清除 pending IRQ。

若需要连续采样，应使用 IIO triggered buffer：硬件 IRQ 或 hrtimer 触发采集，poll function 把带时间戳的扫描数据推入缓冲。用户空间通过标准 sysfs 和 \texttt{/dev/iio:deviceN} 使用驱动，避免形成不可维护的 ioctl ABI。

完整验收包括 raw/scale 换算、参考电压变化、采样超时、并发读取、buffer overflow、runtime suspend、系统休眠、解绑和输入越界。YAML binding 还应通过 \texttt{dt\_binding\_check} 与 \texttt{dtbs\_check}。

\section{并行 NOR Flash 与 MTD 移植}

并行 NOR 的移植不能只向 JEDEC ID 表追加器件编号。应先沿 \texttt{physmap\_flash\_probe()}、\texttt{do\_map\_probe()}、CFI/JEDEC probe 路径确认问题属于器件识别、总线映射还是平台访问语义，再核对总线宽度、字节序、解锁地址、扇区布局和查询时序。

某些平台只允许固定宽度访问，或 CPU 与 Flash 总线字节序不同。这种差异应封装在 \texttt{map\_info} 的 read/write/copy 回调中，而不是散布到通用 MTD 核心。新增识别延时也必须来自数据手册的最大时序要求，不能以某块样板上的偶然延时替代规范依据。

移植完成后至少验证：完整擦除区域表、首尾地址、跨扇区读写、掉电后内容、只读保护、并发访问以及在不同温度和总线频率下的识别稳定性。

\section{电源管理}

系统睡眠处理整机挂起和恢复；runtime PM 在系统运行期间按设备空闲状态关闭资源。
驱动的电源管理回调必须与 I/O、IRQ、DMA 和时钟状态协调，保证挂起时没有新的事务进入，恢复时硬件状态与软件缓存一致。

低功耗问题经常跨越设备树、时钟、pinctrl、regulator、唤醒中断和用户空间策略。
验证时既要检查平均功耗，也要检查唤醒延迟、丢中断、总线错误和重复挂起恢复后的资源泄漏。

\section{BSP 移植流程}

建议按可观测性和依赖关系推进新板移植：
\begin{enumerate}
  \item 确认电源、复位、启动模式和调试串口；
  \item 建立最小 Bootloader 路径并验证 DRAM；
  \item 启动最小内核，保留 early console；
  \item 接入中断控制器、时钟源和系统定时器；
  \item 接入 pinctrl、clock、reset、regulator 等基础供应者；
  \item 接入存储、网络和根文件系统；
  \item 按产品依赖增加显示、音频、传感器等外设；
  \item 完成 runtime PM、系统休眠、压力和恢复测试；
  \item 固化 defconfig、设备树 schema 检查和自动化板级测试。
\end{enumerate}

每一步都应保留已知良好的镜像和日志。出现回归时，优先缩小到电源、时钟、中断、DMA、内存或协议中的一个边界，而不是同时修改多个子系统。

\section{调试与验收}

驱动调试可结合动态调试、tracepoint、ftrace、debugfs、设备子系统统计和总线分析仪。
生产接口不能依赖 debugfs，因为它不属于稳定用户空间 ABI。

BSP 的最低验收条件包括：
\begin{itemize}
  \item 所有设备树均通过 schema 检查且启动日志无未解释错误；
  \item 驱动绑定、解绑和失败回滚不存在资源泄漏或释放后使用；
  \item DMA、缓存、IOMMU 和大内存压力场景经过验证；
  \item 中断风暴、设备超时和热插拔路径能够恢复；
  \item runtime PM、suspend/resume 和唤醒源经过重复测试；
  \item BSP 能由锁定的工具链和配置重新生成相同版本的镜像；
  \item 补丁具有清晰的上游状态和长期维护责任。
\end{itemize}

第~\ref{chap:case-t113-programmer}~章进一步讨论专用 Linux 设备中的资源冲突处理、多实例字符驱动、稳定 ioctl ABI、Buildroot 组装和目标板测试。

第~\ref{chap:embedded-network-time}~章把 MAC/PHY、NAPI、工业网络和硬件时间戳扩展为完整网络数据路径。
